% =====================================================================
%  Single-Pass gem-bcc -- two floats, subroutines in prose.
%  Preamble: \usepackage[ruled,vlined,linesnumbered]{algorithm2e}
%            \usepackage{amsmath,amssymb}
% =====================================================================
\newcommand{\nb}{n_b}
\newcommand{\ns}{n_s}
\newcommand{\pare}{\mathit{par}}
\SetKwFunction{FDem}{de-merge}
\SetKwFunction{FComp}{compress-graph}
\SetKwFunction{FFast}{fast-bcc$^{+}$}
\SetKwFunction{FSF}{Spanning-Forest}
\SetKwFunction{FET}{Euler-Tour}
\SetKwFunction{FTags}{Compute-SF-Tags}
\SetKwFunction{FCC}{CC-Labels}
\SetKwFunction{FPlain}{PlainEdges}
\SetKwFunction{FFen}{Fenced}
\SetKwFunction{FBlk}{is-block}
\newcommand{\cls}{\mathit{cls}}

% ---------------------------------------------------------------- Alg 1
\begin{algorithm}[t]
\caption{Single-Pass \textsc{gem-bcc} Algorithm}\label{alg:sp}
\SetKwInOut{Input}{Input}\SetKwInOut{Output}{Output}
\Input{A graph $G(V,E)$ in edge stream format, batch size $b$}
\Output{BCC labels for edges and cut vertex status for vertices in $G$.}
\lForEach{vertex $v$ \textup{do in parallel}}{$M[v] = v$}
$H_0 \gets (V,\emptyset)$, \quad $\ns \gets |V|$, \quad $\nb \gets 0$
  \tcp*{no spanning tree of $G$ is built}
\For{$i = 1$ \KwTo $k$}{
  $H_{i-1} \gets H_{i-1} + \FDem{$B_i, M, \nb, \ns$}$\;
  $(H_i, I, \nb, \ns) \gets \FComp{$H_{i-1}, \nb$}$\;
  \lForEach{vertex $u \in V(G)$ \textup{do in parallel}}{$M[u] \gets I[M[u]]$}
}
\lForEach{vertex $v \in V(G)$ \textup{do in parallel}}{$\mathit{label}[v] \gets M[v]$}
  \tcp*{round $k$ already left $l$ in the dense ids}
$\mathit{ch}[\cdot] \gets$ root each tree of the forest $H_k$\;
\lForEach{distinct label $j$ \textup{do in parallel}}{
  $\mathit{componentHead}[j] \gets M^{-1}(\mathit{ch}[j])$}
\end{algorithm}

% ---------------------------------------------------------------- Alg 2
\begin{algorithm}[t]
\caption{\textsc{de-merge} \textsc{Algorithm}}\label{alg:dem}
\SetKwInOut{Input}{Input}\SetKwInOut{Output}{Output}
\Input{A batch $B$, mapping $M$, block count $\nb$, next free id $\ns$}
\Output{The edges $B$ contributes to $H$; $M$ and $\ns$ updated}
$E' \gets \emptyset$\;
\ForEach{$(u,v) \in B$ \textup{with} $M[u] \neq M[v]$ \textup{do in parallel}}{
  \ForEach{$x \in \{u,v\}$ \textup{with} $M[x] < \nb$}{
    $L[x] \gets$ a fresh id $\ge \ns$ \tcp*{$x$ leaves its block}
    $E' \gets E' \cup \{(L[x],\,M[x])\}$ \tcp*{de-merge link}
  }
  $E' \gets E' \cup \{(L[u],\,L[v])\}$ \tcp*{the batch edge}
}
$M \gets L$, \quad $\ns \gets \ns + |\{x : x \text{ released}\}|$\;
\BlankLine
\textbf{where}\quad $L[x] = M[x]$ for every $x$ that was not released, and the
fresh ids are allotted by a prefix sum over the released vertices.
\end{algorithm}

% ---------------------------------------------------------------- Alg 3
\begin{algorithm}[t]
\caption{\textsc{compress-graph} \textsc{Algorithm}}\label{alg:comp}
\SetKwInOut{Input}{Input}\SetKwInOut{Output}{Output}
\Input{A forest $H$, block count $\nb$ (ids $<\nb$ are blocks)}
\Output{compressed $H'$, renaming $I$, new counts $\nb'$, $\ns'$}
$l[\cdot] \gets \FFast{$H, \nb$}$ \tcp*{one class per edge of $H$}
\lForEach{node $x$ \textup{do in parallel}}{
  $\mathit{cut}[x] \gets \bigl|\{\,l[e] : e \text{ incident to } x\,\}\bigr| \ge 2$}
\lForEach{class $j$ \textup{do in parallel}}{$\mathit{blk}[j] \gets \FBlk{$j$}$}
\lForEach{node $x$ \textup{do in parallel}}{
  $\mathit{own}[x] \gets \mathit{cut}[x] \vee \neg\,\mathit{blk}[\cls(x)]$
  \tcp*{$x$ survives}}
$\nb' \gets |\{j : \mathit{blk}[j]\}|$, \quad
  $\ns' \gets \nb' + |\{x : \mathit{own}[x]\}|$\;
$I \gets$ prefix sums sending blocks to $[0,\nb')$ and own nodes to
  $[\nb',\ns')$\;
$H' \gets \{\,(I[u],I[v]) : (u,v) \in E(H),\ I[u] \neq I[v]\,\}$\;
\BlankLine
\textbf{where}\quad $\FBlk{$j$}$ holds when class $j$ carries at least two
edges, \emph{or} meets a node $<\nb$, \emph{or} is itself a node $<\nb$ that no
edge of this round touched; and $\cls(x)$ is the class of $x$'s incident edges,
with $\mathit{own}[x]$ set to $[\,x \ge \nb\,]$ when $x$ has none.
\end{algorithm}

% ---------------------------------------------------------------- Alg 4
\begin{algorithm}[t]
\caption{\textsc{fast-bcc}$^{+}$ \textsc{Algorithm}}\label{alg:fast}
\SetKwInOut{Input}{Input}\SetKwInOut{Output}{Output}
\Input{A graph $H$, block count $\nb$ (ids $<\nb$ are blocks)}
\Output{labels $l[\cdot]$ and $\mathit{componentHead}[\cdot]$}
$(E_T, R) \gets \FSF{$H$}$ \tcp*{$H$ is a forest}
$(T,\mathit{first}[\cdot],\mathit{last}[\cdot]) \gets \FET{$E_T, R$}$\;
$(\mathit{low}[\cdot],\mathit{high}[\cdot],F_c) \gets \FTags{$H$}$\;
$G' \gets \FPlain{$T,\nb$}\; \cup\; F_c$ \tcp*{Skeleton graph}
$l[\cdot] \gets \FCC{$G'$}$ \tcp*{Labels for vertices}
\ForEach{vertex $v \in V(H) \setminus R$ \textup{do in parallel}}{
  \lIf{$l[v] \neq l[\pare[v]]$}{$\mathit{componentHead}[l[v]] \gets \pare[v]$}
}
\BlankLine
\textbf{where}\quad
$\FPlain{$T,\nb$} = \bigl\{\,(\pare(v),v) \in T \;:\;
   \underline{\pare(v) < \nb}\;\ \vee\ \neg\,\FFen{$\pare(v),v$}\,\bigr\}$
\end{algorithm}

% ---------------------------------------------------------------- prose
\subsection{Subroutines}

\noindent
\emph{Why \textsc{de-merge} is needed.} \textsc{gem-bcc} builds a spanning tree
of $G$ in a first pass for one reason: if $G_{i-1}$ is disconnected, a non-cut
vertex of $G_{i-1}$ can become a cut vertex of $G_i$, and seeding $H_0$ with a
spanning tree keeps every $G_i$ connected so that this never happens. Dropping
that pass puts the failure back: our $H_0$ is edgeless, so $H_i$ is a forest
throughout and a vertex already absorbed into a block may be separated by a
later batch. \textsc{de-merge} (Alg.~\ref{alg:dem}) is the local repair that
replaces the global pass. It releases from its block every endpoint of the
batch that a new edge can no longer leave inside, giving it a fresh id above
the watermark and a \emph{de-merge link} back to the block it left, so the
block keeps its identity and its remaining residents while the released vertex
is free to join a different component. A block is therefore never split, only
reduced. The computational depth of this subroutine is $O(\log|B|)$, with a
total work complexity of $O(|B|)$.

\medskip\noindent
\emph{Why the dense layout.} \textsc{compress-graph} (Alg.~\ref{alg:comp})
merges the non-cut vertices of each block of $H$ into one node and keeps the
cut vertices as they are. The three clauses of \FBlk{} are what make a block
stable across rounds: an ordinary block holds at least two edges, but a block
reduced by \textsc{de-merge} to a single attachment, or left with no edge at
all, is still a finished bcc and must not be demoted to an ordinary node. The
renumbering is the other half of the contract --- placing blocks in
$[0,\nb')$ and everything else above is what lets the single comparison
$x < \nb$ stand for ``$x$ is a block'' in \textsc{de-merge} and in
\FPlain{}, with no auxiliary flag array. Its depth is $O(\log^2 n)$.

\medskip\noindent
\emph{Reading a class off an edge.} \textsc{fast-bcc}$^{+}$ returns $l[\cdot]$
indexed by vertex, because each non-root vertex owns exactly one parent edge.
For an arbitrary edge $e$ we write $l[e]$ for $l$ at the \emph{deeper} endpoint
of $e$, the one with the larger $\mathit{first}$. On a tree edge
$(\pare(v),v)$ this is $l[v]$, so the two readings agree: $l$ is already an
edge labelling, and \textsc{Compress-Graph} needs no separate array.

\medskip\noindent
\emph{Depth, work and passes.} Each edge of $G$ is copied from cpu to gpu
exactly \emph{once}, so the i/o complexity is $O(m/b)$ -- one pass, against
two for \textsc{gem-bcc} and \textsc{pea-bcc}. No subroutine of the batch loop performs a
bfs, so no term proportional to $\mathit{diameter}(H_i)$ enters: the depth of
bcg construction is $O(k\log^2 n)$ and its work is $O((kn+m)\log n)$. Only the
final head readout traverses $H_k$, once, outside the loop. For each $1 \le i \le
k$, $|E(H_i)| \le |V(H_i)| \le n$ and $|B_i|$ is $O(n)$, so the space
requirement in gpu is $O(n)$.
